%! AP Statistics — Unit 8, Lesson 8.1 — Homework
\documentclass[10pt]{article}
\usepackage{apstats-article}
\usepackage{apstats-boxes}
\usepackage{grading}
\gradingassignment{apstats-u08-l01-hw}
\gradingcourse{advanced_statistics}

\begin{document}

\pageheader{Unit 8, Lesson 8.1}{Homework}

\vspace{0.1in}
\textbf{Directions:} For each problem, show your work and explain your reasoning clearly.

\begin{enumerate}

  \item \textbf{Colored Candies.} A bag of 200 colored candies is supposed to contain
        20\% of each of 5 colors. The observed counts are:

        \vspace{0.05in}
        \renewcommand{\arraystretch}{1.3}
        \begin{tabularx}{0.75\linewidth}{Xccccc}
            \textbf{Color}    & Red & Orange & Yellow & Green & Blue \\
            \hline
            \textbf{Observed} & 45  & 38     & 42     & 31    & 44   \\
        \end{tabularx}

        \begin{enumerate}[label=\alph*.]
          \item Find the expected count for each color. Show your calculation.
                \gradedhide{\writelines{2}\gradedright{\textit{each:}~\gradedblank[id=hw1a, type=short, points=1, answer=40, check=numeric, tolerance=0]{0.9in}{0.2in}}}

          \item Which color shows the largest absolute difference
                $|\text{Observed} - \text{Expected}|$ from expected?
                \gradedhide{\writelines{2}\gradedright{\gradedblank[id=hw1b, type=short, points=1, answer=Green, check=text]{1.2in}{0.2in}}}

          \boxguard[15]
          \item Explain in your own words what question a chi-square goodness-of-fit
                test would answer about this data.
                \gradedhide{\writeline\\\gradedhere{\gradedblank[id=hw1c, type=rubric, points=2, display={Are the gaps from 40 per color too large to be chance, so the bag does not follow 20\% each?}]{\linewidth+6pt}{\gradedlinesheight{2}}}\writelines{9}}
        \end{enumerate}

  \item \textbf{Rolling a Die.} A fair six-sided die is rolled 120 times. You expect
        each face to appear 20 times. The observed counts are:

        \vspace{0.05in}
        \renewcommand{\arraystretch}{1.3}
        \begin{tabularx}{0.75\linewidth}{Xcccccc}
            \textbf{Face}     & 1  & 2  & 3  & 4  & 5  & 6  \\
            \hline
            \textbf{Observed} & 18 & 22 & 19 & 25 & 16 & 20 \\
        \end{tabularx}

        \begin{enumerate}[label=\alph*.]
          \boxguard[15]
          \item Compute the difference (Observed $-$ Expected) for each face.
                \gradedhide{\writeline\\\gradedhere{\gradedblank[id=hw2a, type=rubric, points=2, display={$-2,\ +2,\ -1,\ +5,\ -4,\ 0$}]{\linewidth+6pt}{\gradedlinesheight{1}}}\writelines{9}}

          \boxguard[15]
          \item Would you be surprised if a fair die produced these results?
                Explain without any formal calculation.
                \gradedhide{\writeline\\\gradedhere{\gradedblank[id=hw2b, type=rubric, points=2, display={Not surprising: all gaps are small next to 20 (largest $+5$); chance variation explains them.}]{\linewidth+6pt}{\gradedlinesheight{2}}}\writelines{9}}
        \end{enumerate}

\end{enumerate}

\begin{reflectionbox}
\textbf{Reflection:} In your own words, describe the general question that chi-square
tests answer about categorical data.

\writelines{3}
\end{reflectionbox}

\end{document}
